% ============================================================
%  ICVGIP Submission — ACM sigconf format (acmart)
%  MACE: Multimodal Affective Companion with Empathy
%
%  Converted from the IEEEtran camera-ready version (paper.tex).
%  All experimental results are unchanged and remain verified
%  from executed code.
%
%  BUILD:
%    pdflatex paper_icvgip
%    bibtex   paper_icvgip
%    pdflatex paper_icvgip
%    pdflatex paper_icvgip
%
%  Preamble verified line-by-line against the official
%  ICVGIP-Latex-Template.tex (ICVGIP 2026) shipped in this folder.
%
%  ICVGIP 2026: double-blind; main manuscript limited to 8 pages
%  including text, figures and tables (references may exceed).
% ============================================================
\documentclass[sigconf, review=true]{acmart}
%Do not remove the review=true option for papers submitted for review to ICVGIP2026.

% acmart already loads amsmath, amsfonts, amssymb, booktabs,
% graphicx, xcolor and hyperref. Only additional packages here.
\usepackage{algorithm}
\usepackage{algorithmic}
\usepackage{multirow}
\usepackage{array}

\hyphenation{GoEmotions RoBERTa}

% -----------------------------------------------------------
% Rights / conference metadata — copied verbatim from the
% official ICVGIP-Latex-Template.tex.
% -----------------------------------------------------------
\setcopyright{rightsretained}

% DOI  - Required only for Camera Ready
%\acmDOI{10.475/123_4}

% ISBN - Required only for Camera Ready
%\acmISBN{123-4567-24-567/08/06}

%Conference
\acmConference[ICVGIP2026]{17th Indian Conference on Computer Vision, Graphics and Image Processing}{December 21-25, 2026}{Kolkata, India}
\acmYear{2026}
\copyrightyear{2026}
\acmPrice{15.00}

\begin{document}

% -----------------------------------------------------------
\title[MACE: A Multimodal Affective Companion]{MACE: A Multimodal
Affective Companion with Chain-of-Thought Prompting and Preference
Data Infrastructure for Future RLHF Training}

% -----------------------------------------------------------
% ANONYMISED FOR DOUBLE-BLIND REVIEW.
% ICVGIP 2026 requires that authors not include names,
% affiliations, acknowledgments, or any identifying information.
% UNCOMMENT AND COMPLETE ONLY FOR THE CAMERA-READY VERSION:
%
% \author{Meghana Kuruva}
% \affiliation{%
%   \institution{Indian Institute of Information Technology Sri City}
%   \city{Sri City}
%   \state{Andhra Pradesh}
%   \country{India}}
% \email{meghana.k23@iiits.in}
%
% \author{Piyush Joshi}
% \affiliation{%
%   \institution{Indian Institute of Information Technology Sri City}
%   \city{Sri City}
%   \state{Andhra Pradesh}
%   \country{India}}
% \email{piyush.j@iiits.in}
% -----------------------------------------------------------

\renewcommand{\shortauthors}{}

% -----------------------------------------------------------
\begin{abstract}
Affective computing systems that perceive and respond to human
emotion across multiple modalities hold promise for mental health
support and human--computer interaction, yet most deployed
dialogue systems rely on a single input channel and cannot adapt
to individual user preferences without retraining.
This paper presents MACE (Multimodal Affective Companion with
Empathy), an AI companion that fuses 28-class text emotion
recognition, facial affect inference, and voice interaction within
a Chain-of-Thought (CoT) prompting architecture.
Text emotion is classified using a fine-tuned RoBERTa model on
the GoEmotions corpus; its public model card reports
F1~$\approx$~0.54 under per-label optimised decision thresholds.
Facial affect is inferred zero-shot via Gemini-1.5-Flash using a
constrained single-label prompt---a paradigm that avoids
dataset-specific training while leveraging VLM visual reasoning.
We characterise this design choice against published FER2013
benchmarks (Section~\ref{sec:facial}).
Controlled text classification experiments on an 8-class corpus
($N$=34,792) show that class-weighted LinearSVC achieves the
best macro F1 (60.42\%).
Holding the classifier family fixed, class weighting yields a
clean gain on disgust detection of $+$9.01~pp (26.39\%~$\to$~35.40\%
F1, Logistic Regression); a larger gain of $+$12.21~pp is observed
when additionally switching from an RBF-kernel SVM to a linear
SVC, a comparison confounded by the simultaneous change in kernel
and regularisation.
SMOTE oversampling underperforms class-weighting (57.68\%
vs.\ 60.42\% macro F1) on sparse Bag-of-Words features.
An automated ablation study demonstrates that enriching
linguistic context improves majority-class detection
(fear~$+$2.63~pp, joy~$+$2.11~pp with bigrams) but degrades
minority-class detection (disgust~$-$1.90~pp), motivating MACE's
explicit CoT injection over learned n-gram representations.
All quantitative text-classification results are single-split
estimates without cross-validation; effects below approximately
2~pp should be interpreted cautiously.
MACE additionally collects binary preference labels as
infrastructure for future reward model training and
Direct Preference Optimisation (DPO).
\end{abstract}

%
% The code below should be generated by the tool at
% http://dl.acm.org/ccs.cfm
% Please copy and paste the code instead of the example below.
% The official ICVGIP template leaves this block commented out;
% generate the authoritative XML before the camera-ready version.
% The concept IDs below are indicative and NOT verified against
% the ACM CCS tool.
%
% \begin{CCSXML}
% <ccs2012>
%    <concept>
%        <concept_id>10010147.10010178.10010179</concept_id>
%        <concept_desc>Computing methodologies~Natural language processing</concept_desc>
%        <concept_significance>500</concept_significance>
%        </concept>
%    <concept>
%        <concept_id>10010147.10010178.10010224</concept_id>
%        <concept_desc>Computing methodologies~Computer vision</concept_desc>
%        <concept_significance>300</concept_significance>
%        </concept>
%    <concept>
%        <concept_id>10010147.10010257</concept_id>
%        <concept_desc>Computing methodologies~Machine learning</concept_desc>
%        <concept_significance>300</concept_significance>
%        </concept>
%    <concept>
%        <concept_id>10003120.10003121</concept_id>
%        <concept_desc>Human-centered computing~Human computer interaction (HCI)</concept_desc>
%        <concept_significance>300</concept_significance>
%        </concept>
%  </ccs2012>
% \end{CCSXML}
%
% \ccsdesc[500]{Computing methodologies~Natural language processing}
% \ccsdesc[300]{Computing methodologies~Computer vision}
% \ccsdesc[300]{Computing methodologies~Machine learning}
% \ccsdesc[300]{Human-centered computing~Human computer interaction (HCI)}

\keywords{affective computing, emotion recognition, multimodal
fusion, chain-of-thought prompting, reinforcement learning from
human feedback, large language models, FER2013, GoEmotions, class
imbalance, vision-language models}

\maketitle

% -----------------------------------------------------------
\section{Introduction}
\label{sec:intro}

Human emotion is a complex, multidimensional phenomenon that
manifests simultaneously across linguistic, paralinguistic, and
physiological channels~\cite{ekman1992}.
Emotional intelligence in conversational agents---the capacity to
detect affective states and respond with appropriate
empathy---remains an open challenge at the intersection of natural
language processing (NLP), affective computing, and
human-computer interaction (HCI).

Large language models (LLMs)~\cite{devlin2019bert,liu2019roberta}
have enabled richer conversational systems, yet most deployed
architectures treat emotion as a coarse binary (positive/negative)
or a three-class label.
Such granularity discards psychologically meaningful distinctions
between, for example, \emph{grief} and \emph{disappointment}, or
\emph{admiration} and \emph{joy}, that are essential for genuine
empathetic engagement.
Furthermore, existing affective dialogue systems are predominantly
unimodal and static: they process one input channel and cannot
adapt response style to individual preferences without full model
retraining.

This work addresses these limitations through five contributions:

\begin{enumerate}
    \item \textbf{Fine-grained multimodal emotion perception.}
          MACE classifies 28 fine-grained emotions from text
          using a fine-tuned RoBERTa model~\cite{samlow2022goemotions}
          and infers facial affect via a vision-language model
          (VLM) using zero-shot constrained prompting.
    \item \textbf{Facial inference benchmarking.}
          We characterise VLM-based facial emotion inference
          against the FER2013 benchmark~\cite{goodfellow2013fer},
          contrasting the zero-shot VLM paradigm with published
          supervised CNN and transformer baselines.
    \item \textbf{Chain-of-Thought contextual prompting.}
          A structured CoT template integrates text emotion,
          facial expression, user memory, and conversation history
          to condition an LLM for empathetic response
          generation~\cite{wei2022chainofthought}.
    \item \textbf{Class-imbalance-aware classification.}
          Controlled experiments demonstrate that class-weighted
          loss outperforms SMOTE on sparse text features,
          recovering 9--12 F1 points on disgust depending on the
          classifier family; the larger figure reflects an
          additional kernel/solver change alongside weighting
          (Section~\ref{sec:results}).
    \item \textbf{Preference data infrastructure for future RLHF.}
          Binary helpfulness signals are stored persistently,
          forming a preference corpus for future reward model
          training via DPO~\cite{ouyang2022rlhf,rafailov2023dpo}.
\end{enumerate}

Section~\ref{sec:related} reviews related work.
Section~\ref{sec:methodology} describes the system architecture
and mathematical formulation.
Section~\ref{sec:algorithm} presents the inference algorithm.
Section~\ref{sec:setup} details the experimental setup.
Section~\ref{sec:results} reports text classification results.
Section~\ref{sec:facial} presents facial inference benchmarking.
Section~\ref{sec:ablation} presents the ablation study.
Section~\ref{sec:comparison} provides comparative analysis.
Section~\ref{sec:userstudy} describes the planned user evaluation.
Section~\ref{sec:limitations} acknowledges limitations, and
Section~\ref{sec:conclusion} concludes.

% -----------------------------------------------------------
\section{Related Work}
\label{sec:related}

\subsection{Text-Based Emotion Recognition}

Early approaches relied on lexical resources such as the NRC
Emotion Lexicon~\cite{mohammad2013crowdsourcing} and rule-based
systems grounded in Ekman's six basic
emotions~\cite{ekman1992}.
The SemEval-2007 Affective Text task~\cite{strapparava2007semeval}
stimulated supervised learning with SVM classifiers.
Pre-trained transformers---BERT~\cite{devlin2019bert} and
RoBERTa~\cite{liu2019roberta}---subsequently demonstrated that
contextual embeddings transfer effectively to emotion classification.
Demszky et al.~\cite{demszky2020goemotions} introduced
GoEmotions: 58,000 Reddit comments annotated across 28 fine-grained
emotion labels, enabling substantially richer taxonomies and
reporting an average F1 of $\approx$0.46 for their fine-tuned
BERT model on the full 28-class task.

\subsection{Facial Emotion Recognition}

Goodfellow et al.~\cite{goodfellow2013fer} introduced FER2013,
establishing that human accuracy on posed facial expression images
is only $\approx$65.5\%, reflecting the intrinsic difficulty of
the task.
Subsequent CNN architectures progressively improved accuracy from
$\approx$65\% to $\approx$73\%~\cite{li2022fer_survey}.
Transformer-based approaches, including Vision Transformer (ViT)
variants, have pushed state-of-the-art above
85\%~\cite{li2022fer_survey}.
Conventional approaches are trained in a supervised manner
per-dataset; by contrast, vision-language
models~\cite{alayrac2022flamingo} enable zero-shot facial
emotion inference by leveraging broad pre-trained visual
understanding, trading potential accuracy ceiling for
deployment flexibility and the elimination of dataset-specific
training.

\subsection{Class Imbalance in Affective Corpora}

Affective datasets are inherently imbalanced: high-frequency
states such as joy co-occur far more commonly than rare states
such as disgust or shame.
Class-weighted loss and SMOTE~\cite{chawla2002smote} oversampling
are standard mitigation strategies.
However, SMOTE's effectiveness degrades on high-dimensional sparse
feature spaces~\cite{blagus2013smote} because interpolating
between nearest neighbours in BoW space yields semantically
incoherent documents.

\subsection{Multimodal Affective Dialogue}

Poria et al.~\cite{poria2016convolutional} demonstrated that
fusing textual, acoustic, and visual features outperforms any
single modality for sentiment analysis.
Hazarika et al.~\cite{hazarika2018conversational} introduced
memory networks for conversational emotion recognition, showing
that discourse history materially benefits multimodal encoders.

\subsection{Chain-of-Thought Prompting and RLHF}

Wei et al.~\cite{wei2022chainofthought} showed that prompting
LLMs with explicit intermediate reasoning steps substantially
improves multi-step task performance.
Ouyang et al.~\cite{ouyang2022rlhf} demonstrated that
Reinforcement Learning from Human Feedback (RLHF)---through
reward modelling and Proximal Policy Optimisation (PPO)---aligns
LLM outputs with human preferences more effectively than
supervised fine-tuning alone.
Rafailov et al.~\cite{rafailov2023dpo} unified reward modelling
and policy optimisation into Direct Preference Optimisation (DPO),
a single-stage contrastive objective that is the target of
MACE's future training pipeline.

\subsection{Positioning}

None of the systems above combine fine-grained multimodal emotion
perception, benchmark-characterised zero-shot facial inference,
persistent cross-session memory, and a preference-logging
infrastructure within a single deployed conversational agent.
MACE does not claim to advance the state of the art on any
individual sub-task (text emotion classification, facial emotion
recognition, or CoT prompting); its contribution is the
integration and the controlled analysis of the design trade-offs
this integration exposes, detailed in Section~\ref{sec:comparison}.

% -----------------------------------------------------------
\section{System Methodology}
\label{sec:methodology}

\subsection{Architecture Overview}

MACE operates as a five-stage pipeline:
\textbf{(1)} multimodal input acquisition (text, optional image,
optional voice transcription);
\textbf{(2)} emotion perception (text: 28-class RoBERTa;
face: zero-shot VLM);
\textbf{(3)} context retrieval (conversation history, user memory,
preference exemplars from a local SQLite database);
\textbf{(4)} CoT prompt construction; and
\textbf{(5)} LLM response generation with binary preference
feedback capture.
Algorithm~\ref{alg:mace} formalises this pipeline for a single
conversational turn.

\subsection{Text Emotion Classification}

Let $x \in \mathcal{X}$ denote a natural language utterance.
The text emotion classifier maps $x$ to a probability
distribution over $K{=}28$ GoEmotions categories via fine-tuned
RoBERTa:
\begin{equation}
    \mathbf{p}^{\text{text}} = \text{softmax}\!\left(
    \mathbf{W} \cdot \text{RoBERTa}_\theta(x) + \mathbf{b}
    \right),
    \label{eq:text_emotion}
\end{equation}
where $\mathbf{W} \in \mathbb{R}^{K \times d}$ are learned
projection weights and $d$ is the encoder hidden dimension.
The primary detected emotion and its confidence score are:
\begin{equation}
    e^* = \arg\max_{k} \, p^{\text{text}}_k, \quad
    c^* = \max_{k} \, p^{\text{text}}_k.
    \label{eq:primary_emotion}
\end{equation}

\subsection{Facial Emotion Inference}

Given an optional camera image $I$, a VLM $g_\phi$ returns a
single-label facial affect estimate under a constrained
prompt $\mathcal{P}_{\text{face}}$:
\begin{equation}
    e^{\text{face}} = g_\phi\!\left(I,\; \mathcal{P}_{\text{face}}\right).
    \label{eq:face_emotion}
\end{equation}
When no image is available, $e^{\text{face}}{=}\varnothing$ and
this field is omitted from the downstream CoT prompt.
The prompt $\mathcal{P}_{\text{face}}$ constrains the output to a
single emotion label from a fixed vocabulary, preventing
free-form responses that would require post-hoc parsing.

\subsection{Persistent User Memory}

Let $\mathcal{T}{=}\{(\phi_j, c_j)\}$ denote trigger
phrase--category pairs, $c_j \in \{\texttt{like},
\texttt{dislike}\}$.
A preference item is extracted and stored as:
\begin{equation}
    m_i = \text{parse}(x, \phi_j), \quad
    \mathcal{M} \leftarrow \mathcal{M} \cup \{(m_i, c_j)\},
    \label{eq:memory}
\end{equation}
where $\text{parse}(\cdot)$ extracts the noun phrase following
trigger $\phi_j$.
$\mathcal{M}$ is persisted across sessions in SQLite with
insertion-time deduplication.

\subsection{Chain-of-Thought Prompt Construction}

Let $\mathcal{H}_{t-1}$ denote the three most recent
$(u_i, r_i)$ conversation turns, and let
$\mathcal{F}^+, \mathcal{F}^-$ denote the five most recent
positively- and negatively-rated interaction pairs retrieved
from the preference database.
The CoT context vector at turn $t$ is:
\begin{equation}
    \Phi_t = \bigl(e^*,\; c^*,\; e^{\text{face}},
               \; \mathcal{M},\; \mathcal{H}_{t-1},\;
               \mathcal{F}^+,\; \mathcal{F}^-\bigr).
    \label{eq:cot_context}
\end{equation}
The response is generated by conditioning a pre-trained LLM
$\mathcal{G}$ on the serialised context:
\begin{equation}
    r_t = \mathcal{G}\!\left(\, \pi(\Phi_t),\; x_t \right),
    \label{eq:generation}
\end{equation}
where $\pi(\cdot)$ serialises $\Phi_t$ into a structured
natural-language prompt that exposes the reasoning chain
explicitly before generating the empathetic reply.

\subsection{Preference Data Collection}

After each response $r_t$, the user provides a binary signal
$y_t \in \{+1, -1\}$ (``helpful'' / ``not helpful'').
The tuple $(x_t, \Phi_t, r_t, y_t)$ is appended to the
preference database $\mathcal{D}$.

\textbf{Current deployment scope.}
MACE does not yet train a reward model or update $\mathcal{G}$
via gradient descent.
Positive preference exemplars are injected into the CoT prompt
as in-context demonstrations, providing an immediate alignment
signal.
This constitutes test-time preference adaptation rather than
full RLHF.
The formal DPO training objective over preference
pairs $(\mathbf{y}_w, \mathbf{y}_l)$ that will be applied
once $|\mathcal{D}|$ is sufficient is:
\begin{equation}
    \mathcal{L}_{\text{DPO}} =
    -\mathbb{E}\!\left[\log\sigma\!\left(
    \beta \log \frac{\mathcal{G}(\mathbf{y}_w \mid x)}
                    {\mathcal{G}_0(\mathbf{y}_w \mid x)}
    - \beta \log \frac{\mathcal{G}(\mathbf{y}_l \mid x)}
                      {\mathcal{G}_0(\mathbf{y}_l \mid x)}
    \right)\right],
    \label{eq:dpo}
\end{equation}
where $\mathcal{G}_0$ is the frozen reference policy and
$\beta{>}0$ is a KL-penalty coefficient~\cite{rafailov2023dpo}.

% -----------------------------------------------------------
\section{Algorithm}
\label{sec:algorithm}

Algorithm~\ref{alg:mace} summarises the MACE inference pipeline
for a single conversational turn.

\begin{algorithm}
\caption{MACE: Single-Turn Inference}
\label{alg:mace}
\begin{algorithmic}[1]
\REQUIRE Utterance $x_t$; optional image $I_t$;
         preference database $\mathcal{D}$; memory $\mathcal{M}$
\ENSURE  Response $r_t$; updated $\mathcal{D}$, $\mathcal{M}$

\STATE $\mathbf{p}^{\text{text}} \leftarrow
       f_\theta(x_t)$ \COMMENT{Eq.~\eqref{eq:text_emotion}}
\STATE $(e^*, c^*) \leftarrow
       (\arg\max_k p^{\text{text}}_k,\;
        \max_k p^{\text{text}}_k)$
\IF{$I_t \neq \varnothing$}
    \STATE $e^{\text{face}} \leftarrow
           g_\phi(I_t, \mathcal{P}_{\text{face}})$
\ELSE
    \STATE $e^{\text{face}} \leftarrow \varnothing$
\ENDIF
\STATE $(m_{\text{new}}, c_{\text{new}}) \leftarrow
       \text{MemoryExtract}(x_t)$ \COMMENT{Eq.~\eqref{eq:memory}}
\IF{$m_{\text{new}} \neq \varnothing$}
    \STATE $\mathcal{M} \leftarrow \mathcal{M} \cup
           \{(m_{\text{new}}, c_{\text{new}})\}$
\ENDIF
\STATE $\mathcal{H}_{t-1} \leftarrow
       \text{RetrieveHistory}(\mathcal{D},\; n{=}3)$
\STATE $(\mathcal{F}^+, \mathcal{F}^-) \leftarrow
       \text{RetrievePreferences}(\mathcal{D},\; k{=}5)$
\STATE $\Phi_t \leftarrow (e^*, c^*, e^{\text{face}},
       \mathcal{M}, \mathcal{H}_{t-1}, \mathcal{F}^+,
       \mathcal{F}^-)$ \COMMENT{Eq.~\eqref{eq:cot_context}}
\STATE $r_t \leftarrow \mathcal{G}(\pi(\Phi_t),\; x_t)$
       \COMMENT{Eq.~\eqref{eq:generation}}
\STATE $\text{SaveChat}(\mathcal{D},\; x_t, r_t, e^*, c^*)$
\STATE $y_t \leftarrow \text{CollectBinaryFeedback}()$
\STATE $\text{SavePreference}(\mathcal{D},\; y_t)$
\RETURN $r_t$
\end{algorithmic}
\end{algorithm}

% -----------------------------------------------------------
\section{Experimental Setup}
\label{sec:setup}

\subsection{Text Emotion Dataset}

Baseline text classification experiments use an 8-class emotion
corpus of $N{=}34{,}792$ utterances.
Table~\ref{tab:class_dist} shows the class distribution, which
is severely imbalanced: \textit{joy} accounts for 31.75\% of
samples while \textit{shame} contributes only 0.42\%.
Given this imbalance, macro-averaged F1 is the primary evaluation
metric; accuracy is reported solely for completeness.

\begin{table}[htbp]
\centering
\caption{Class Distribution of the 8-Class Text Corpus}
\label{tab:class_dist}
\begin{tabular}{lrr}
\toprule
\textbf{Emotion} & \textbf{Samples} & \textbf{Proportion (\%)} \\
\midrule
Joy      & 11{,}045 & 31.75 \\
Sadness  &  6{,}722 & 19.32 \\
Fear     &  5{,}410 & 15.55 \\
Anger    &  4{,}297 & 12.35 \\
Surprise &  4{,}062 & 11.68 \\
Neutral  &  2{,}254 &  6.48 \\
Disgust  &    856   &  2.46 \\
Shame    &    146   &  0.42 \\
\midrule
\textbf{Total} & \textbf{34{,}792} & \textbf{100.00} \\
\bottomrule
\end{tabular}
\end{table}

Raw utterances are preprocessed using NeatText~\cite{neattext2021}:
\textit{(i)} user handles (\texttt{@mentions}) and
\textit{(ii)} English stop-words are removed sequentially.
A stratified 70/30 split (\texttt{random\_state=42}) yields
24,353 training and 10,438 test samples.

\subsection{Text Classifiers}

Seven experimental conditions are evaluated, all using
CountVectorizer as the text representation layer.

\begin{itemize}
    \item \textbf{LR}: Logistic Regression, L-BFGS solver,
          $C{=}1.0$, \texttt{max\_iter=1000}.
    \item \textbf{LR-bal}: LR with \texttt{class\_weight=`balanced'}.
    \item \textbf{SVM}: SVM, RBF kernel, $C{=}10$.
    \item \textbf{SVC-bal}: Linear SVC,
          \texttt{class\_weight=`balanced'},
          \texttt{max\_iter=2000}.
    \item \textbf{RF}: Random Forest~\cite{breiman2001rf},
          100 estimators, \texttt{random\_state=42} (a common
          default ensemble size for tabular/text baselines).
    \item \textbf{RF-bal}: RF with
          \texttt{class\_weight=`balanced'}.
    \item \textbf{SMOTE+LR}: SMOTE~\cite{chawla2002smote}
          applied in feature space, followed by LR.
\end{itemize}

\subsection{Deployed Text Emotion Model}

MACE's production classifier is
\texttt{SamLowe/roberta-base-go\_emotions}~\cite{samlow2022goemotions},
a RoBERTa-base model fine-tuned on the GoEmotions
corpus~\cite{demszky2020goemotions} (58,000 utterances,
28 emotion classes).
The model card reports F1~$\approx$~0.45 at a uniform 0.5
decision threshold, rising to F1~$\approx$~0.54 when each label's
threshold is individually optimised; we did not independently
re-evaluate this model and report the card's own figures for
transparency.
The model is loaded via the HuggingFace \texttt{transformers}
pipeline with \texttt{top\_k=None} to obtain the full
probability distribution over all 28 classes.

\textbf{Comparability note.}
The 8-class BoW baseline experiments and the 28-class RoBERTa
results are on distinct datasets and label spaces and cannot be
compared numerically.
The BoW baselines demonstrate the ceiling imposed by
non-contextual representations, motivating the transformer-based
production model.

\subsection{FER2013 Facial Emotion Dataset}
\label{sec:fer_setup}

The Facial Expression Recognition 2013 (FER2013)
dataset~\cite{goodfellow2013fer} was introduced at the ICML 2013
Representation Learning Workshop.
It comprises 35,887 greyscale 48$\times$48-pixel facial images
collected from web sources, annotated with one of seven emotion
categories: Angry, Disgust, Fear, Happy, Neutral, Sad, and
Surprise.
The official partition provides 28,709 training, 3,589
validation, and 3,589 test images.
We use the standard test split for all benchmark comparisons.

% -----------------------------------------------------------
\section{Text Emotion Classification Results}
\label{sec:results}

\subsection{Overall Performance}

Table~\ref{tab:summary} reports accuracy, macro F1, and
weighted F1 for all seven conditions.
SVC-bal achieves the best macro F1 (60.42\%) and is
recommended as the primary BoW baseline.

\begin{table}[htbp]
\centering
\caption{Text Classifier Performance ($n_{\text{test}}{=}10{,}438$)}
\label{tab:summary}
\begin{tabular}{lrrr}
\toprule
\textbf{Model} & \textbf{Acc (\%)} & \textbf{Mac-F1 (\%)} &
\textbf{Wt-F1 (\%)} \\
\midrule
LR          & 63.54 & 59.48 & 62.75 \\
LR-bal      & 61.10 & 60.06 & 61.44 \\
SVM         & 62.96 & 58.44 & 61.90 \\
SVC-bal     & 61.36 & \textbf{60.42} & 61.27 \\
RF          & 58.32 & 55.12 & 57.97 \\
RF-bal      & 53.68 & 54.77 & 55.17 \\
SMOTE+LR    & 59.91 & 57.68 & 60.23 \\
\bottomrule
\end{tabular}
\end{table}

\subsection{Per-Class F1: Effect of Class Weighting}

Table~\ref{tab:perclass} contrasts standard and class-weighted
variants for the LR and SVC families.

\begin{table*}[t]
\centering
\caption{Per-Class F1 (\%): Standard vs.\ Class-Weighted Classifiers}
\label{tab:perclass}
\begin{tabular}{lrrrrrr}
\toprule
\textbf{Emotion} &
\textbf{LR} & \textbf{LR-bal} & \textbf{$\Delta$} &
\textbf{SVM} & \textbf{SVC-bal} & \textbf{$\Delta$} \\
\midrule
Joy      & 69.47 & 67.84 & $-$1.63 & 68.46 & 67.70 & $-$0.76 \\
Sadness  & 58.57 & 57.18 & $-$1.39 & 57.60 & 56.87 & $-$0.73 \\
Fear     & 70.93 & 70.75 & $-$0.18 & 71.82 & 70.60 & $-$1.22 \\
Anger    & 59.01 & 58.21 & $-$0.80 & 56.11 & 57.02 & $+$0.91 \\
Surprise & 49.50 & 49.75 & $+$0.25 & 48.26 & 48.33 & $+$0.07 \\
Neutral  & 66.62 & 55.99 & $-$10.63& 68.98 & 60.78 & $-$8.20 \\
Disgust  & 26.39 & \textbf{35.40} & $+$9.01  &
           20.98 & \textbf{33.19} & $+$12.21 \\
Shame    & 75.32 & 85.39 & $+$10.07 & 75.32 & 88.89 & $+$13.57 \\
\midrule
\textbf{Macro F1} & 59.48 & 60.06 & $+$0.58 &
                    58.44 & \textbf{60.42} & $+$1.98 \\
\bottomrule
\end{tabular}
\end{table*}

\textbf{Class weighting recovers minority classes; the SVM-family
comparison is confounded.}
Within the Logistic Regression family, LR and LR-bal differ only
in \texttt{class\_weight}, making this pair a clean, single-variable
ablation: disgust F1 rises from 26.39\% to 35.40\% ($+$9.01~pp)
and shame from 75.32\% to 85.39\% ($+$10.07~pp).
This is the primary evidence in this paper that class weighting
increases the gradient contribution from minority-class errors
during optimisation, counteracting the frequency-driven bias of
unweighted training.
The SVM-family pair (SVM~$\to$~SVC-bal) shows a larger apparent
gain on disgust (20.98\%~$\to$~33.19\%, $+$12.21~pp) and shame
(75.32\%~$\to$~88.89\%, $+$13.57~pp); however, this comparison
changes three factors simultaneously---kernel (RBF~$\to$~linear),
the regularisation constant ($C{=}10$~$\to$~default $C{=}1$), and
class weighting---so we do not attribute the full SVM-family
delta to weighting alone.
We report both pairs for completeness, but treat the LR-family
result as the controlled estimate of the weighting effect and
the SVM-family result as a comparison between two
deployment-plausible configurations rather than an ablation.
Class weighting also comes at a cost to neutral classification:
balanced LR drops neutral F1 from 66.62\% to 55.99\%
($-$10.63 pp).
In a therapeutic companion context this trade-off is arguably
acceptable---neutral misclassification is less harmful than
silently failing to detect disgust or shame---though we have not
tested this assumption with human raters.

\textbf{SMOTE underperforms class-weighting.}
SMOTE+LR achieves 57.68\% macro F1, below both LR-bal
(60.06\%) and SVC-bal (60.42\%).
SMOTE synthesises minority-class examples via interpolation
between $k$-nearest neighbours in feature space; on sparse
high-dimensional BoW vectors this produces semantically
incoherent token distributions~\cite{blagus2013smote},
which adds noise rather than signal.

\textbf{Random Forest degrades under balancing.}
RF-bal (53.68\% accuracy, 54.77\% macro F1) underperforms
unbalanced RF (58.32\%, 55.12\%).
In high-dimensional sparse spaces, each tree selects features
randomly; balanced weights amplify minority-class gradients
but simultaneously increase split-level variance, degrading
majority-class recall disproportionately.

\subsection{Qualitative System Interaction Examples}

Table~\ref{tab:qual} presents illustrative MACE responses
showing how the CoT context modulates response generation.

\begin{table*}[t]
\centering
\caption{Illustrative MACE Interaction Examples (Single- and Multi-Turn)}
\label{tab:qual}
\begin{tabular}{p{2.6cm}p{2.1cm}p{1.8cm}p{1.6cm}p{5.7cm}}
\toprule
\textbf{User Utterance} & \textbf{Text Emotion (Conf.)} &
\textbf{Facial} & \textbf{Memory} & \textbf{System Response} \\
\midrule
``I feel so alone tonight.'' &
sadness (0.91) & Sad & None &
``I hear you, and I want you to know your feelings are completely
 valid. Feeling alone can be incredibly heavy---I am right
 here with you, and you do not have to carry this alone.'' \\
\addlinespace
``I love basketball---I played so well today!'' &
joy (0.87) & Happy & basketball (like) &
``That is wonderful! Your passion for basketball really shines
 through---enjoy this feeling, you earned it.
 \emph{(Noted: you like basketball.)}'' \\
\midrule
\multicolumn{5}{l}{\textit{Multi-turn example (conversation
history $\mathcal{H}_{t-1}$ active):}}\\[2pt]
Turn 1: ``I hate my job.'' &
anger (0.79) & Neutral & None &
``That sounds really frustrating. Would you like to talk about
 what specifically has been getting to you?'' \\
Turn 2: ``My manager never listens.'' &
anger (0.74) & Neutral &
\mbox{Hist: anger} &
``It makes complete sense that you are angry---feeling unheard
 by someone in authority is genuinely demoralising.
 Your frustration from earlier is still coming through.'' \\
\bottomrule
\end{tabular}
\end{table*}

The multi-turn example illustrates the direct value of
$\mathcal{H}_{t-1}$: without history, the response to Turn~2
would treat the complaint as isolated.
With history active, the system references prior emotional
continuity explicitly (``Your frustration from earlier''),
a qualitative improvement not captured by the automated
proxy metrics in Section~\ref{sec:ablation}.

% -----------------------------------------------------------
\section{Facial Emotion Inference Validation}
\label{sec:facial}

\subsection{Motivation}

MACE's facial inference module employs Gemini-1.5-Flash via
a constrained zero-shot prompt rather than a supervised
facial emotion recognition (FER) model.
This section contextualises this design choice by:
\textit{(i)} reporting published supervised performance on
FER2013 as the benchmark landscape;
\textit{(ii)} comparing the methodological trade-offs between
zero-shot VLM and supervised inference; and
\textit{(iii)} analysing the failure modes known to affect
zero-shot VLM inference on facial images.
This paper does not report a quantitative accuracy figure for
Gemini-1.5-Flash on FER2013 or any other benchmark.

\subsection{Supervised Baseline Performance on FER2013}

Table~\ref{tab:fer_baselines} reports published accuracy on
the FER2013 test set for representative supervised methods
spanning the CNN-to-transformer evolution, alongside the
human performance benchmark.

\begin{table}[htbp]
\centering
\caption{Published Supervised Performance on FER2013 Test Set}
\label{tab:fer_baselines}
\begin{tabular}{llr}
\toprule
\textbf{Method} & \textbf{Type} & \textbf{Acc (\%)} \\
\midrule
Human~\cite{goodfellow2013fer}            & Human          & 65.5 \\
Deep CNN~\cite{goodfellow2013fer}         & CNN (2013)     & 65.4 \\
VGGNet fine-tuned~\cite{li2022fer_survey} & CNN transfer   & 70.1 \\
ResNet variants~\cite{li2022fer_survey}   & CNN transfer   & 72.4 \\
Transformer (ViT)~\cite{li2022fer_survey} & ViT fine-tuned & 87.8 \\
\bottomrule
\end{tabular}
\end{table}

Two observations from Table~\ref{tab:fer_baselines} merit
emphasis.
First, human accuracy (65.5\%) is \emph{not} substantially
above early CNN performance (65.4\%), underscoring the
inherent ambiguity of posed facial expressions.
Second, transformer-based approaches require full supervised
fine-tuning on FER2013 training data to achieve 87.8\%
accuracy; this training phase is entirely absent in
zero-shot VLM inference.

\subsection{VLM vs.\ Supervised FER: Methodological Comparison}

Zero-shot VLM inference and supervised FER models represent
two fundamentally different operating regimes, each with
distinct trade-offs:

\begin{itemize}
    \item \textbf{Data independence.}
          Supervised models require hundreds to tens of
          thousands of labelled facial images per target domain.
          A VLM requires no domain-specific training data,
          enabling immediate deployment in novel environments
          (e.g., clinical settings, new demographic groups).
    \item \textbf{Label space flexibility.}
          Changing the emotion taxonomy for a supervised model
          requires full retraining.
          For the VLM approach, the taxonomy is modified by
          editing $\mathcal{P}_{\text{face}}$.
    \item \textbf{Accuracy ceiling.}
          Supervised transformers trained on FER2013 data
          achieve $\approx$88\% accuracy by exploiting
          task-specific training signal unavailable to
          zero-shot models.
          Zero-shot VLMs sacrifice this ceiling for generality.
    \item \textbf{Resolution sensitivity.}
          FER2013 images are 48$\times$48 pixels, which is
          below the resolution at which modern VLMs are most
          reliable for fine-grained visual tasks.
          Real-world webcam frames are typically higher
          resolution but may introduce additional confounders.
\end{itemize}

In the MACE deployment context, the VLM approach is preferred
because:
\textit{(i)} MACE is not trained on any FER dataset and
must generalise to arbitrary users and environments;
\textit{(ii)} the emotion label vocabulary can be adapted
to match the same 28-class GoEmotions taxonomy as the text
classifier by modifying $\mathcal{P}_{\text{face}}$;
and \textit{(iii)} the facial signal is auxiliary---it
enriches the CoT prompt but is not the primary emotion
source---so moderate inference accuracy is acceptable.

\subsection{Failure Mode Analysis}

Based on published analyses of VLM performance on facial
emotion tasks~\cite{li2022fer_survey} and the known properties
of FER2013, the following failure modes are anticipated:

\textbf{Fear--Surprise confusion.}
Fear and surprise share raised eyebrows and widened eyes;
these are visually similar cues that confuse both human
annotators and classifiers.
Goodfellow et al.~\cite{goodfellow2013fer} report this as
the dominant confusion pair on FER2013.

\textbf{Anger--Disgust confusion.}
Both emotions produce similar brow-lowering and
lip-compression expressions.
Zero-shot models lacking exposure to fine-grained FAU
labelling are particularly susceptible to this confusion.

\textbf{Neutral--Sadness confusion.}
Subtle sadness expressions may be scored as neutral by
models lacking fine-grained training.
This is especially prevalent for low-intensity emotional
expressions, which are common in naturalistic webcam
captures.

\textbf{Resolution degradation.}
FER2013's 48$\times$48 pixel resolution renders fine-grained
cues (e.g., subtle lip curvature, sclera exposure) largely
invisible.
While webcam frames are typically higher resolution,
unfavourable lighting or motion blur can similarly obscure
these cues.

\textbf{Pose and occlusion.}
FER2013 images predominantly display frontal, unoccluded
faces.
Real-world interactions involve off-axis head poses, partial
occlusion by glasses or hair, and varied lighting.
Zero-shot VLMs may generalise better to these conditions
than supervised models trained predominantly on frontal
posed images; empirical verification on in-the-wild datasets
(e.g., AffectNet~\cite{mollahosseini2019affectnet}) constitutes
future work.

% -----------------------------------------------------------
\section{Ablation Study}
\label{sec:ablation}

\subsection{Motivation and Design}

A full system ablation---disabling each CoT component and
rating responses with human annotators---requires Gemini API
calls and recruited participants; this is reserved as part of
the planned user study (Section~\ref{sec:userstudy}).
We present an automated proxy ablation that is fully executable
and captures the central research question: \emph{does
enriching the linguistic context available to a classifier
improve detection accuracy, and which emotion classes benefit
most?}

We operationalise \emph{context richness} via n-gram scope,
using LR-bal throughout to isolate context effects from
classifier choice (SVC-bal achieves marginally higher macro F1
but does not alter the directional findings):
\begin{itemize}
    \item \textbf{V1 (Unigram)}: text-only with no contextual
          word-order dependencies.
    \item \textbf{V2 (Unigram+Bigram)}: captures adjacent-word
          relationships; serves as a controlled proxy for the
          effect of incorporating short-range textual context.
    \item \textbf{V3 (Unigram+Bigram+Trigram)}: fullest
          surface-level context.
\end{itemize}

\subsection{Ablation Results}

Table~\ref{tab:ablation} reports overall metrics and
Table~\ref{tab:ablation_perclass} shows the per-class
transition from V1 to V2.

\begin{table}[htbp]
\centering
\caption{Ablation: N-Gram Context Depth (LR-bal)}
\label{tab:ablation}
\begin{tabular}{lccc}
\toprule
\textbf{Variant} & \textbf{N-gram} & \textbf{Acc (\%)} &
\textbf{Mac-F1 (\%)} \\
\midrule
V1 --- Unigram [text-only]          & (1,1) & 61.10 & 60.06 \\
V2 --- +Bigram [+context proxy]     & (1,2) & \textbf{62.78} &
\textbf{60.39} \\
V3 --- +Trigram [full surface]      & (1,3) & 62.64 & 59.64 \\
\bottomrule
\end{tabular}
\end{table}

\begin{table}[htbp]
\centering
\caption{Per-Class F1 (\%) Change: V1$\to$V2 (LR-bal)}
\label{tab:ablation_perclass}
\begin{tabular}{lrrl}
\toprule
\textbf{Emotion} & \textbf{V1 (\%)} & \textbf{V2 (\%)} &
\textbf{$\Delta$} \\
\midrule
Fear     & 70.75 & 73.38 & $+$2.63 $\uparrow$ \\
Joy      & 67.84 & 69.95 & $+$2.11 $\uparrow$ \\
Surprise & 49.75 & 51.28 & $+$1.53 $\uparrow$ \\
Sadness  & 57.18 & 58.53 & $+$1.35 $\uparrow$ \\
Anger    & 58.21 & 58.85 & $+$0.64 $\uparrow$ \\
Shame    & 85.39 & 83.72 & $-$1.67 $\downarrow$ \\
Neutral  & 55.99 & 53.95 & $-$2.04 $\downarrow$ \\
Disgust  & 35.40 & 33.50 & $-$1.90 $\downarrow$ \\
\midrule
Macro F1 & 60.06 & 60.39 & $+$0.33 \\
\bottomrule
\end{tabular}
\end{table}

\subsection{Discussion and Architectural Implications}

\textbf{Majority classes benefit from richer context.}
Adding bigrams improves F1 for fear ($+$2.63~pp), joy ($+$2.11~pp),
and surprise ($+$1.53~pp).
These classes have $n \geq 4{,}000$ training samples, providing
sufficient coverage to learn informative bigram patterns
(e.g., ``not happy'', ``so scared'', ``can't believe'').

\textbf{Minority classes degrade with richer context.}
Disgust F1 drops 1.90 pp from V1 to V2, and neutral drops
2.04 pp.
With only $\approx$599 disgust training examples, the bigram
feature space is extremely sparse: most co-occurrence patterns
appear once or twice, introducing noise.
Trigrams exacerbate this: disgust drops further to 29.12\%
at V3 versus 35.40\% at V1.

\textbf{Architectural implication.}
Learned n-gram representations disproportionately harm
minority classes when context scope is widened, because
those classes lack sufficient data to populate multi-token
feature patterns.
MACE's CoT architecture directly addresses this:
rather than \emph{learning} context-dependent associations
from labelled data, it \emph{explicitly provides} the
detected emotion, facial signal, and history as structured
natural language to the LLM.
This eliminates the data-sparsity bottleneck for minority
classes.
The ablation provides quantitative motivation for preferring
explicit contextual injection over learned n-gram features.

\textbf{V3 exhibits diminishing returns.}
Moving from V2 to V3 does not improve macro F1
(60.39\%~$\to$~59.64\%), consistent with high-order n-gram
behaviour on short emotional texts where most trigrams are
hapax legomena.

% -----------------------------------------------------------
\section{Comparative Analysis}
\label{sec:comparison}

Table~\ref{tab:comparison} positions MACE against representative
affective dialogue systems on six architectural dimensions.

\begin{table*}[t]
\centering
\caption{Architectural Comparison of Affective Dialogue Systems}
\label{tab:comparison}
\begin{tabular}{lcccccc}
\toprule
\textbf{System} &
\textbf{Modalities} &
\textbf{Emotion Classes} &
\textbf{Persistent Memory} &
\textbf{Preference Feedback} &
\textbf{LLM Response} &
\textbf{FER Validation} \\
\midrule
SemEval-2007~\cite{strapparava2007semeval} &
Text & 6 & No & No & No & No \\
Poria et al.~\cite{poria2016convolutional} &
T+A+V & 7 & No & No & No & Yes \\
Hazarika et al.~\cite{hazarika2018conversational} &
T+A+V & 7 & Episodic & No & No & Yes \\
Standard instruction-tuned LLM &
Text & Implicit & Session only & No & Yes & No \\
\textbf{MACE (proposed)} &
\textbf{T+V+Voice} &
\textbf{28} &
\textbf{Persistent} &
\textbf{Binary (DPO-ready)} &
\textbf{Yes} &
\textbf{No (\S\ref{sec:facial})} \\
\bottomrule
\end{tabular}
\end{table*}

Among the systems surveyed, MACE is the only one combining
fine-grained 28-class text emotion recognition, a zero-shot VLM
facial inference module characterised against (but not yet
quantitatively validated on) FER2013, cross-session persistent
memory, multimodal voice interaction, and a structured preference
data pipeline for future DPO alignment.

% -----------------------------------------------------------
\section{Planned User Evaluation}
\label{sec:userstudy}

\subsection{Preliminary Qualitative Observations}

During development, the authors informally exercised MACE across
a range of hand-crafted single- and multi-turn scenarios spanning
the 8 emotion classes; two representative transcripts are shown
in Table~\ref{tab:qual}.
These development-time sessions were not systematically logged or
scored and therefore do not constitute quantitative evidence; we
report them only as motivation for the structured protocol below,
not as a validated result.
Qualitatively, enabling conversation history caused responses to
reference prior emotional context explicitly (Table~\ref{tab:qual},
multi-turn example), a pattern absent when history was disabled.
A systematic, logged measurement of memory-referencing accuracy
and emotion-acknowledgement rate across a fixed, pre-registered
scenario set is deferred to the IRB-approved study described next.

\subsection{Full Evaluation Protocol}

\textbf{Participants.}
$N{=}30$ adult participants will be recruited following
Institutional Review Board (IRB) approval, with a minimum
of 15 per condition in a within-subjects design.

\textbf{Conditions.}
Each participant interacts with Full MACE and a text-only
baseline (no facial input, no persistent memory, no history)
across 10 standardised emotional vignettes in counterbalanced
order.

\textbf{Primary measure.}
Perceived empathy is assessed via the Jefferson Scale of
Empathy adapted for AI agents~\cite{hojat2001jse}
(7-point Likert, 5 items).

\textbf{Secondary measures.}
Contextual coherence (3-item 5-point Likert);
personalisation recall (binary: did the system correctly
reference a stated preference?); and task completion time.

\textbf{Analysis.}
Paired $t$-tests for between-condition empathy score
differences; McNemar's test for personalisation recall
accuracy.
Results will be reported in a follow-up publication.

% -----------------------------------------------------------
\section{Limitations}
\label{sec:limitations}

\textbf{Facial inference: benchmark context vs.\ deployment.}
Section~\ref{sec:facial} characterises published FER2013
supervised baselines to contextualise the trade-off of using a
zero-shot VLM rather than a supervised FER model; no accuracy or
F1 figure for Gemini-1.5-Flash is claimed anywhere in this paper.
Furthermore, FER2013 consists of controlled, predominantly
frontal greyscale images collected under limited pose and
lighting variation.
Real-world webcam interaction introduces additional
confounders: motion blur, oblique head pose, partial
occlusion, variable illumination, and compressed video
artefacts.
Benchmark performance should therefore be treated as an
upper bound on in-the-wild accuracy.
Evaluation on naturalistic datasets such as
AffectNet~\cite{mollahosseini2019affectnet} constitutes
future work.

\textbf{Text classifier BoW ceiling.}
The BoW baselines plateau at 60.42\% macro F1 for SVC-bal;
the deployed contextual transformer (RoBERTa-base) is reported,
per its model card, at F1~$\approx$~0.45--0.54 (depending on
thresholding) on the substantially harder 28-class GoEmotions
task, underscoring the non-comparability of these evaluation
settings.
Future work should evaluate RoBERTa on the same 8-class corpus
to enable direct, controlled comparison.

\textbf{Automated proxy for ablation.}
The n-gram ablation is verifiable and reproducible but is
a proxy for response empathy, not a direct measurement.
The full component ablation---removing each CoT element and
eliciting human empathy ratings---remains as part of the
planned user study.

\textbf{Single-split evaluation; no significance testing.}
All text-classification results (Tables~\ref{tab:summary}--%
\ref{tab:ablation_perclass}) are computed on one stratified
70/30 split with a fixed seed (\texttt{random\_state=42}), without
cross-validation, repeated resampling, or significance testing.
Differences below approximately 2~pp---for example, the
surprise-class delta between SVM and SVC-bal ($+$0.07~pp) or the
anger-class delta between LR and LR-bal ($-$0.80~pp)---should be
treated as within the likely range of split-to-split variance
rather than as established effects.
The larger effects reported for disgust, shame, and neutral are
less likely to be pure noise given their magnitude, but this has
not been confirmed with repeated runs.
Future work should report means and standard deviations over
multiple seeds or cross-validation folds before treating any
single-split delta below this threshold as conclusive.

\textbf{RLHF pipeline not yet active.}
Binary preference data is collected but has not yet been
used to train a reward model or to fine-tune $\mathcal{G}$.
Full DPO training~\cite{rafailov2023dpo} requires a preference
corpus of at least several thousand labelled pairs.

\textbf{Privacy and data governance.}
User utterances and facial images are transmitted to external
commercial APIs.
Applications involving sensitive mental health disclosures
should instead employ a locally hosted open-weight
LLM~\cite{ouyang2022rlhf}, which would retain all data
on-device.

\textbf{Neutral-class precision--recall trade-off.}
Class weighting reduces neutral F1 by up to 10.63 pp
(Table~\ref{tab:perclass}).
This trade-off is acceptable in the therapeutic companion
context but would require re-evaluation if MACE were applied
to high-neutral domains such as factual question answering.

% -----------------------------------------------------------
\section{Conclusion}
\label{sec:conclusion}

This paper presented MACE, a multimodal empathetic AI companion
that integrates 28-class text emotion recognition via a
fine-tuned RoBERTa model (GoEmotions model-card F1~$\approx$~0.54
under per-label optimised thresholds); zero-shot facial affect
inference via Gemini-1.5-Flash, whose design is motivated by a
characterisation of the published FER2013 benchmark landscape;
voice interaction; Chain-of-Thought contextual prompting with
persistent user memory; and a binary preference data
infrastructure designed for future DPO alignment.

Controlled text classification experiments on an 8-class corpus
($N{=}34{,}792$) established that class-weighted LinearSVC
achieves the best macro F1 (60.42\%).
A clean, single-variable ablation (Logistic Regression with and
without class weighting) shows disgust F1 rising from 26.39\% to
35.40\% ($+$9.01~pp) when class weighting is the only factor
changed; a larger gain of $+$12.21~pp is observed when
additionally switching from an RBF-kernel SVM to a linear SVC,
a comparison confounded by the simultaneous kernel and
regularisation change and reported here as a practical
configuration comparison rather than an isolated ablation.
SMOTE oversampling underperforms class-weighting (57.68\%
macro F1), consistent with reported degradation of synthetic
interpolation on sparse BoW features~\cite{blagus2013smote}.

The automated ablation study quantified a measurable asymmetry:
enriching linguistic context via bigrams improves majority
classes (fear $+$2.63~pp, joy $+$2.11~pp) but degrades minority
classes (disgust $-$1.90 pp), plausibly due to feature sparsity.
This motivates, but does not by itself establish, that MACE's
explicit CoT injection outperforms learned n-gram representations
for empathetic response quality; that comparison requires the
human-rated component ablation described in
Section~\ref{sec:userstudy}.
All text-classification results derive from a single 70/30 split
without cross-validation (Section~\ref{sec:limitations}).

The facial inference section characterises the published FER2013
benchmark landscape---human accuracy of 65.5\%, supervised
CNN accuracy of $\approx$65--72\%, and ViT accuracy of
$\approx$88\%---to motivate the zero-shot VLM design choice on
methodological grounds (data independence, label-space
flexibility) rather than on a quantitative accuracy comparison.
No accuracy or F1 figure for Gemini-1.5-Flash is reported in this
paper.

Future work encompasses: \textit{(i)} the planned 30-participant
IRB-approved user study with Jefferson Scale empathy ratings,
which will also provide the first systematic, logged measurement
of memory-referencing accuracy in place of the informal
observations in Section~\ref{sec:userstudy};
\textit{(ii)} full component ablation with human raters;
\textit{(iii)} DPO training on accumulated preference
data~\cite{rafailov2023dpo};
\textit{(iv)} evaluation of the facial inference module on
AffectNet~\cite{mollahosseini2019affectnet}; and
\textit{(v)} a privacy-preserving on-device LLM deployment.

% -----------------------------------------------------------
\bibliographystyle{ACM-Reference-Format}
\bibliography{paper_icvgip}

\end{document}
